\section{Editorial: Complex Types}
\label{sec:complex-types}

Instead of a primitive like \texttt{string} or \texttt{int}, a field can
declare \texttt{complex: ClassName} --- an escape hatch to a hand-written
type living outside the generated file, imported from a real,
language-native location rather than generated by Emi itself. A built-in
type such as \texttt{Date} can be used directly; a user-defined type like
\texttt{Money} has to be registered in the module's \texttt{complexes:}
list, naming the class and where to import it from.

\begin{lstlisting}
name: emiComplexFields
complexes:
  - name: Money
    location: ../some/directory/Money
actions:
  - name: sample
    method: post
    out:
      fields:
        - name: date
          complex: +Date
        - name: money
          complex: +Money
\end{lstlisting}

Prefixing the complex name with \texttt{+} (\texttt{+Date}, \texttt{+Money})
tells Emi to \textbf{instantiate} the value when parsing a response,
instead of merely typing it. For the JS compiler this produces a real
setter that coerces on assignment:

\begin{lstlisting}
set money(value: Money) {
  if (value instanceof Money) {
    this.#money = value;
  } else {
    this.#money = new Money(value);   // coerced, never trusted blindly
  }
}
\end{lstlisting}

\texttt{Date} is treated as a built-in complex and is never imported;
\texttt{Money} is imported from the given \texttt{location} and
instantiated on assignment. Every generated setter for a complex field
(\texttt{setDate}, \texttt{setMoney}, \ldots) returns \texttt{this}, so
assignments chain.

\begin{pullquote}
On the Go side, a \texttt{complex} field used inside an \texttt{entity}
has to pull its own weight: it must implement
\texttt{database/sql/driver.Valuer} and \texttt{sql.Scanner} itself if it
needs to persist as anything other than gorm's default guess --- the
generator can't know how to store an arbitrary hand-written type
(\S\ref{sec:go-entities}).
\end{pullquote}
